\subsection{Euclidean division of polynomials in one indeterminate}

PROPOSITION 10. --- Let \(f\) and \(g\) be non-zero elements of \(\mathbf{A}[\mathbf{X}]\) of degrees \(m\) and \(n\) respectively. Let \(\alpha_0\) be the leading coefficient of \(f\) and \(\mu = \sup (n - m + 1, 0)\) . There exist \(u, v \in \mathbf{A}[\mathbf{X}]\) such that

\[
\alpha_ {0} ^ {\mu} g = u f + v, \quad \deg v <   m.
\]

If \(\alpha_0\) is cancellable in A, then \(u\) and \(v\) are uniquely determined by these properties.

The existence of \(u\) and \(v\) is clear when \(n < m\) , because then we can take \(u = 0\) and \(v = g\) . For \(n \geqslant m\) we shall use induction on \(n\) . Let \(\beta\) be the leading coefficient of \(g\) ; if \(f = \sum_{k=0}^{m} \alpha_k X^{m-k}\) , we can write \(\alpha_0^\mu g = \alpha_0^{\mu-1} \beta X^{n-m} f + \alpha_0^{\mu-1} g_1\) , where \(g_1 \in A[X]\) and \(\deg g_1 < n\) . By the induction hypothesis there exist \(u_1, v \in A[X]\) such that \(\alpha_0^{\mu-1} g_1 = u_1 f + v\) and \(\deg v < m\) . Hence

\[
\alpha_ {0} ^ {\mu} g = \left(\alpha_ {0} ^ {\mu - 1} \beta X ^ {n - m} + u _ {1}\right) f + v
\]

and it suffices to put \(u = \alpha_0^{\mu - 1}\beta X^{n - m} + u_1\) .

Assuming \(\alpha_0\) to be cancellable in \(\mathbf{A}\) , let us now prove the uniqueness of \(u\) and \(v\) . Let \(u, v, u_1, v_1 \in \mathbf{A}[\mathbf{X}]\) be such that

\[
\alpha_ {0} ^ {\mu} g = u f + v = u _ {1} f + v _ {1}, \quad \deg v <   m, \quad \deg v _ {1} <   m.
\]

We have \((u - u_{1})f = v_{1} - v\) and \(\deg (v_1 - v) < m\) , hence \(u - u_{1} = 0\) (IV, p.~9, Prop. 7) and therefore \(v_{1} - v = 0\) .

COROLLARY (« Euclidean division of polynomials »). --- Let \(f\) be a non-zero element of A{[}X{]} whose leading coefficient is invertible and \(m = \deg f\) .

\begin{enumerate}
\def\labelenumi{(\roman{enumi})}
\tightlist
\item
  For every \(g \in \mathbf{A}[\mathbf{X}]\) there exist \(u, v \in \mathbf{A}[\mathbf{X}]\) such that
\end{enumerate}

\[
g = u f + v, \quad \deg v <   m.
\]

Moreover, these conditions determine \(u\) and \(v\) uniquely.

\begin{enumerate}
\def\labelenumi{(\roman{enumi})}
\setcounter{enumi}{1}
\item
  The sub-A-modules \(A + AX + \cdots + AX^{m-1}\) and \(fA[X]\) of A{[}X{]} are supplementary in A{[}X{]}.
\item
  Assume \(f\) non-constant and consider \(\mathbf{A}[\mathbf{X}]\) as an \(\mathbf{A}[\mathbf{T}]\) -module by means of the homomorphism \(u(\mathbf{T}) \mapsto u(f(\mathbf{X}))\) of \(\mathbf{A}[\mathbf{T}]\) into \(\mathbf{A}[\mathbf{X}]\) . Then \(\mathbf{A}[\mathbf{X}]\) is a free \(\mathbf{A}[\mathbf{T}]\) -module with basis \((1, \mathbf{X}, \ldots, \mathbf{X}^{m-1})\) .
\end{enumerate}

Assertions (i) and (ii) are immediate consequences of Prop. 10.

Let us prove (iii). Let \(\psi\) be the homomorphism \(v \mapsto v(f(X), X)\) of A{[}T, X{]} into A{[}X{]}. Consider A{[}T, X{]} first as polynomial ring in T with coefficients in A{[}X{]}; Cor. 1 of IV, p.~5 shows that the kernel \(\mathfrak{a}\) of \(\psi\) is the ideal (T - f(X)) of A{[}T, X{]}. Consider now A{[}T, X{]} as polynomial ring in X with coefficients in A{[}T{]}; then \(\psi\) is an A{[}T{]}-linear mapping of A{[}T{]}{[}X{]} into A{[}X{]}. Assertion (ii) above (applied to the polynomial \(f(X) - T\) in X with coefficients in A{[}T{]}) shows that \((1, X, ..., X^{m-1})\) is a basis of an A{[}T{]}-submodule of A{[}T, X{]} supplementary to \(\mathfrak{a}\) . Since \(\psi(X^i) = X^i\) for every integer \(i \geqslant 0\) , (iii) follows at once.

With the notations of (i) we shall say that u is the quotient and v the remainder in the Euclidean division of g by f; for the remainder to vanish it is necessary and sufficient that f should divide g.

